\documentclass[10pt,a4paper]{amsart}
\usepackage[margin=19mm]{geometry}
\usepackage{amsmath,amssymb,mathtools}
\usepackage[scr=rsfs,cal=euler]{mathalfa}
\usepackage{xfrac}
\usepackage[unicode,hidelinks,bookmarks=false]{hyperref}
\newcommand{\ie}{i.e.\ }
\newcommand{\cf}{cf.\ }
\newcommand{\resp}{respectively\ }
\newcommand{\Subsection}{\subsection}
\providecommand{\nobreakdash}{\nobreak\hbox{-}}
\setlength{\emergencystretch}{2em}
\multlinegap0pt
\usepackage{bm}
\usepackage{bbm}
\usepackage{dsfont}
\usepackage{xspace}
\usepackage{cancel}
\usepackage{mathtools}
\usepackage{amsthm}
\makeatletter
\@ifundefined{theorem}{\newtheorem{theorem}{Theorem}}{}
\@ifundefined{lemma}{\newtheorem{lemma}[theorem]{Lemma}}{}
\makeatother
\DeclareMathOperator{\Tr}{Tr}
\title[Pseudo-Riemannian Spectral Triples for $\mathrm{SU}(1,1)$]{Pseudo-Riemannian Spectral Triples for $\mathrm{SU}(1,1)$}
\author{Jort de Groot}
\date{Source: arXiv:2601.22171v1 (CC BY 4.0)}
\begin{document}
\maketitle

\noindent\emph{Source and licence.} Pseudo-Riemannian Spectral Triples for $\mathrm{SU}(1,1)$. Version arXiv:2601.22171v1 (CC BY 4.0), distributed under the Creative Commons Attribution 4.0 International licence, as recorded in the arXiv licence metadata for this exact version and shown on the abstract page https://arxiv.org/abs/2601.22171v1. Full rights notice: licenses/arxiv-2601.22171-CC-BY-4.0.txt.\par\smallskip

\noindent\textit{Editorial scope.} Counted unit: the Lemma whose source label is \texttt{compact\_resolvent} in the article (source0.tex lines 642--644, source label \texttt{compact\_resolvent}), delivered together with the article's own complete original proof of it (source0.tex lines 645--680, one \texttt{proof} environment of 36 lines). No mathematics is written, repaired or re-derived: statement and proof are the article's own text line for line. Required context retained uncounted: compactness\_powers (lines 531--533); eq\_action\_R (lines 547--553); series\_1 (lines 570--576); series\_2 (lines 589--595); series\_3 (lines 604--610). Every definition, display and label of those lines is retained unchanged. Omitted from this package and not redistributed: 1--530, 534--546, 554--569, 577--588, 596--603, 611--641, 681--847. Nothing inside a retained excerpt is omitted or altered: every retained body is delivered exactly as the article's lines are written, so no omission receipt of any kind accompanies this package. Citations: the retained excerpts carry no citation command at all, so no bibliography entry of the article is transcribed and no reference is redistributed. Reference adaptation: none was needed and none was made; every reference inside the delivered unit resolves inside it, so no unresolved reference or citation is produced. Printed slips kept literal: no printed word, symbol, hypothesis or number of the counted statement or of its proof was altered. Layout and class: the article's own class is \texttt{amsart} with packages including mathtools, enumitem, amssymb; the wrapper is the lane's pinned \texttt{amsart} editorial wrapper at 10pt on A4, which restates the theorem-like environments the retained text needs and adds the article's own macro definitions that the retained text needs (\texttt{\textbackslash Tr}), with extra wrapper packages (bm, bbm, dsfont, xspace, cancel, mathtools). The delivered numbering is the unit's own. Assets: no retained span contains a figure environment, a graphics-inclusion command, a PDF-page inclusion, a table or a TikZ picture; no image file is copied into this package, the package declares no asset dependency and no image file is redistributed with it. Licence: version arXiv:2601.22171v1 of this article is distributed under the Creative Commons Attribution 4.0 International licence, as recorded in the arXiv licence metadata for this exact version and shown on the abstract page https://arxiv.org/abs/2601.22171v1. A full rights notice is delivered at licenses/arxiv-2601.22171-CC-BY-4.0.txt. Characters: every retained excerpt is pure ASCII, so the delivered engine, XeLaTeX with its default Latin Modern text font, reports no missing character.
\begin{lemma}\label{compactness_powers}
    Let $A$ be a bounded operator on $\mathcal{H}$ and let $B$ be bounded and self-adjoint. If $AB^q$ is compact for some $q\in\mathbb{R}_{>0}$, then it is compact for all $q>0$.
\end{lemma}

\begin{equation}\label{eq_action_R}
    \begin{aligned}
        R_\pm(T^\chi)&=\sum_m\Big(\tfrac{1}{2}|\tau|^2+m'(m'\pm1)+\tfrac{9}{8}\Big)^{-1/2} v^\chi_m\otimes (v^\chi_m)^*,\\
        R_\pm(T^{l,-})&=\sum_m \Big(-\tfrac{1}{2}l(l+1)+m'(m'\pm1)+\tfrac{9}{8}\Big)^{-1/2} v_m^{l,-}\otimes (v_m^{l,-})^*,\\
        R_\pm(T^{l,+})&=\sum_m \Big(-\tfrac{1}{2}l(l+1)+m'(m'\pm1)+\tfrac{9}{8}\Big)^{-1/2} v_m^{l,+}\otimes (v_m^{l,+})^*.
    \end{aligned}
\end{equation}

\begin{lemma}\label{series_1}
    The series
    \begin{equation*}
        S^-_\pm\coloneq\sum_{l\leq-\tfrac{1}{2}}\Big(-l-\tfrac{1}{2}\Big)\sum_m \frac{1}{\Big(-\tfrac{1}{2}l(l+1)+m'(m'\pm 1)+\tfrac{9}{8}\Big)^{q/2}}
    \end{equation*}
    with $l\in\tfrac{1}{2}\mathbb{Z}_{<0}$, $m<l-\epsilon$, $m'=m+\epsilon$ and $\epsilon=0$ when $l\in\mathbb{Z}$, while $\epsilon=\tfrac{1}{2}$ when $l\in\tfrac{1}{2}\mathbb{Z}_\mathrm{odd}$, converges for $q>3$. 
\end{lemma}

\begin{lemma}\label{series_2}
    The series
    \begin{equation*}
        S^+_\pm\coloneq\sum_{l\leq-\tfrac{1}{2}}\Big(-l-\tfrac{1}{2}\Big)\sum_m\frac{1}{\Big(-\tfrac{1}{2}l(l+1)+m'(m'\pm 1)+\tfrac{9}{8}\Big)^{q/2}}
    \end{equation*}
    with $m\geq-l-\epsilon$ and $l,m',\epsilon$ as in Lemma \ref{series_1}, converges for $q>3$.
\end{lemma}

\begin{lemma}\label{series_3}
    The series
    \begin{equation*}
        S^\chi_\pm=\sum_{\epsilon=0,\tfrac{1}{2}}\int_0^\infty \sum_m\frac{1}{\Big(\tfrac{1}{2}\rho^2+m'(m'\pm1)+\tfrac{5}{4}\Big)^{q/2}}\rho\tanh\pi(\rho+i\epsilon)d\rho
    \end{equation*}
    with $m\in\mathbb{Z}$ and $m'=m+\epsilon$ converges for $q>3$.
\end{lemma}

\begin{lemma}\label{compact_resolvent}
    Let $a\in\mathcal{A}$ and $R=(\langle D\rangle^2+1)^{-1/2}$ the resolvent of $\langle D\rangle$. Then $aR$ is a compact operator on $\mathcal{H}$.
\end{lemma}

\begin{proof}
    We show that $aR^q$ is Hilbert--Schmidt for $q>3$, which using Lemma \ref{compactness_powers} implies compactness of $aR$. Write $a=f\otimes 1$ for $f\in C_c^\infty(\mathrm{SU}(1,1))$; then $aR^q$ is Hilbert--Schmidt if both $fR_+^q$ and $fR_-^q$ are Hilbert--Schmidt operators on $L^2(\mathrm{SU}(1,1))$. We show that $fR_+^q$ is Hilbert--Schmidt for $q>3$. 

    Let $\varphi\in\mathcal{C}(\mathrm{SU}(1,1))$, then $\hat{\varphi}\in\mathcal{C}(\widehat{\mathrm{SU}(1,1)})$ is trace class, so it has a well-defined inverse Fourier transform. Using \eqref{eq_action_R}, the action of $fR_+^q$ is
    \begin{equation*}
        fR_+^q\varphi(g)=\int f(g)\Tr\bigg[R_+^q(\pi)\pi(g)\int\varphi(h)\pi(h^{-1})dh\bigg]d\mu(\pi).
    \end{equation*}
    To write $fR_+^q$ as an integral operator, we need to swap the order of integration. Note that $R_+^q$ is trace class for $q>1$, as is easily checked from \eqref{eq_action_R}, so that we can rewrite
    \begin{equation*}
        fR_+^q\varphi(g)=\int f(g) \int\Tr\big[R_+^q(\pi)\pi(gh^{-1})\big]\varphi(h)dhd\mu(\pi).
    \end{equation*}
    To swap the two integrals, we apply Fubini's theorem. As $g$ is fixed, we consider
    \begin{equation*}
        \int\int\Big|\Tr\big[R_+^q(\pi)\pi(gh^{-1})\big]\varphi(h)\Big|dhd\mu(\pi)\leq \int\Tr\big[R_+^q(\pi)\big]d\mu(\pi)\int|\varphi(h)|dh
    \end{equation*}
    Writing out the trace and the Plancherel measure gives
    \begin{equation*}
        \int\Tr\big[R_+^q(\pi)\big]d\mu(\pi)=S^+_++S^-_++S^\chi_+,
    \end{equation*}
    which converges by Lemmas \ref{series_1}, \ref{series_2}, and \ref{series_3}, as we assumed $q>3$. Hence, by Fubini's theorem, we obtain
    \begin{equation*}
        fR_+^q\varphi(g)=\int f(g)\int\Tr\big[R_+^q(\pi)\pi(gh^{-1})\big]d\mu(\pi)\varphi(h)dh\eqcolon\int K(g,h)\varphi(h)dh.
    \end{equation*}

    From Lemmas \ref{series_1}-\ref{series_3}, it follows that $R_+^q\in L^2(\widehat{\mathrm{SU}(1,1)})$. Hence, using the Plancherel theorem, $\int\Tr\big[R_+^q(\pi)\pi(gh^{-1})\big]d\mu(\pi)$ is an element of $L^2(\mathrm{SU}(1,1))$ (in the variable $h$). As $f\in L^2(\mathrm{SU}(1,1))$, we conclude that $K_+(g,h)\in L^2(\mathrm{SU}(1,1)\times \mathrm{SU}(1,1))$.
    
    For $\psi\in L^2(\mathrm{SU}(1,1))$, let $(\varphi_m)\subset\mathcal{C}(\mathrm{SU}(1,1))$ be a sequence converging to $\psi$ in the $L^2$-norm, then also $fR_+^q \varphi_m\rightarrow fR_+^q\psi$ in $L^2(\mathrm{SU}(1,1))$. There is a subsequence, also denoted $\varphi_m$, that converges pointwise almost everywhere. Using that $K_+$ is an $L^2$-kernel, we find
    \begin{equation*}
        \lim_{m\rightarrow\infty}(fR_+^q\varphi_m)(g)=\lim_{m\rightarrow\infty}\int K_+(g,h)\varphi_m(h)dh=\int K_+(g,h)\psi(h)dh,
    \end{equation*}
    which we use to write
    \begin{equation*}
        fR_+^q\psi(g)=\int K_+(g,h)\psi(h)dh.
    \end{equation*}
    Thus, $fR_+^q$ is an integral operator with kernel in $L^2(\mathrm{SU}(1,1)\times \mathrm{SU}(1,1))$, so that $fR_+^q$ is Hilbert--Schmidt. As the proof for $fR_-^q$ is similar, this concludes the proof.
\end{proof}

\end{document}
